\documentclass[11pt]{article}

\usepackage[a4paper,margin=1in]{geometry}
\usepackage{amsmath}
\usepackage{amssymb}
\usepackage{amsthm}
\usepackage{array}
\usepackage{booktabs}

\newtheorem{theorem}{Theorem}
\newtheorem{lemma}[theorem]{Lemma}
\newtheorem{proposition}[theorem]{Proposition}
\newtheorem{corollary}[theorem]{Corollary}
\theoremstyle{definition}
\newtheorem{definition}[theorem]{Definition}
\newtheorem{remark}[theorem]{Remark}

\newcommand{\Cul}{C}
\newcommand{\Woo}{W}
\newcommand{\Z}{\mathbb{Z}}
\DeclareMathOperator{\lpf}{lpf}

\title{Disproof of conjecture \texttt{00000000226}\\
\large Cullen and Woodall numbers do not have complementary proportions\\
with least prime factor \(3\)}
\author{earthking11}
\date{2026-09-13}

\begin{document}
\maketitle

\begin{abstract}
The conjecture asserts that the proportion of Cullen numbers \(C_n=n2^n+1\)
with least prime factor \(3\) is an explicit rational, and that the
corresponding proportions for the Cullen and Woodall numbers
\(\Woo_n=n2^n-1\) sum to \(1\) (``complementary symmetry mod \(3\)'').
We refute the second claim unconditionally and elementarily. Both sequences
are odd, so ``least prime factor \(3\)'' is equivalent to ``divisible by
\(3\)''. Reducing modulo \(3\) with period \(6\), we prove
\[
3\mid C_n \iff n\equiv 1,2 \pmod 6,
\qquad
3\mid \Woo_n \iff n\equiv 4,5 \pmod 6 .
\]
Each proportion is therefore \(2/6=1/3\), and their sum is
\(1/3+1/3=2/3\neq 1\). The asserted complementary symmetry fails. The result
is formalised in core Lean~4 and verified numerically for all \(n\le 2000\).
\end{abstract}

\section{The conjecture}

\begin{definition}
For \(n\ge 1\), the \emph{Cullen numbers} and \emph{Woodall numbers} are
\[
\Cul_n = n\cdot 2^n + 1,
\qquad
\Woo_n = n\cdot 2^n - 1 .
\]
\end{definition}

The conjecture, as filed in \texttt{conjectures/00000000226.md}, states
(translating the definition and the second claim):

\begin{quote}
\itshape
The count of Cullen primes is asymptotically \(x/\log^2 x\) times an explicit
sieve constant; the proportion of Cullen numbers with least prime factor \(3\)
is an explicit rational, and the corresponding proportions for Cullen and
Woodall numbers sum to \(1\) (complementary symmetry mod \(3\)).
\end{quote}

We do not address the asymptotic prime-counting claim. We refute the claim
about the proportions: the Cullen proportion is exactly \(1/3\), the Woodall
proportion is exactly \(1/3\), and \(1/3+1/3=2/3\neq 1\). Hence the asserted
complementary symmetry modulo \(3\) is false.

\section{Parity: ``least prime factor \(3\)'' means ``divisible by \(3\)''}

\begin{lemma}\label{lem:odd}
For every \(n\ge 1\), \(\Cul_n\) and \(\Woo_n\) are odd.
\end{lemma}

\begin{proof}
Since \(n\ge 1\), the term \(n\cdot 2^n\) is divisible by \(2\) and hence
even. Adding \(1\) gives an odd number, and subtracting \(1\) from an even
number gives an odd number:
\[
\Cul_n = \underbrace{n2^n}_{\text{even}} + 1 \equiv 1 \pmod 2,
\qquad
\Woo_n = \underbrace{n2^n}_{\text{even}} - 1 \equiv 1 \pmod 2. \qedhere
\]
\end{proof}

\begin{corollary}\label{cor:lpf}
For \(n\ge 1\), the least prime factor of \(\Cul_n\) equals \(3\) if and only
if \(3\mid \Cul_n\). The same holds for \(\Woo_n\).
\end{corollary}

\begin{proof}
By Lemma~\ref{lem:odd} the number \(N\in\{\Cul_n,\Woo_n\}\) is odd, so \(2\) is
not a prime factor of \(N\). If \(3\mid N\), then among the prime factors of
\(N\) the smallest is the smallest prime divisor, and since \(2\) is excluded
and \(3\) divides \(N\), that smallest divisor is \(3\). Conversely, if the
least prime factor is \(3\), then \(3\mid N\) by definition.
\end{proof}

Thus, for both sequences, counting numbers with least prime factor \(3\) is the
same as counting numbers divisible by \(3\).

\section{Reduction modulo \(3\) and the residue table}

We work modulo \(3\). Since \(2\equiv -1\pmod 3\),
\begin{equation}\label{eq:twopow}
2^n \equiv (-1)^n \pmod 3
=
\begin{cases}
1, & n \text{ even},\\
2, & n \text{ odd}.
\end{cases}
\end{equation}
The map \(n\mapsto (n \bmod 3,\, n \bmod 2)\) has combined period
\(\operatorname{lcm}(3,2)=6\). Equivalently, \(n\bmod 6\) determines both
\(n\bmod 3\) and the parity of \(n\). Therefore both sequences are periodic
modulo \(3\) with period dividing \(6\), and the truth value of
\(3\mid \Cul_n\) and of \(3\mid \Woo_n\) depends only on \(n\bmod 6\).

The following table records, for each residue class \(n\equiv r\pmod 6\), the
value of \(2^n\bmod 3\), of \(\Cul_n\bmod 3\), and of \(\Woo_n\bmod 3\).

\begin{center}
\begin{tabular}{c c c c}
\toprule
\(r = n\bmod 6\) & \(2^n \bmod 3\) & \(\Cul_n \bmod 3\) & \(\Woo_n \bmod 3\) \\
\midrule
0 & 1 & \(0\cdot 1+1 = 1\) & \(0\cdot 1-1 = -1\equiv 2\) \\
1 & 2 & \(1\cdot 2+1 = 3\equiv 0\) & \(1\cdot 2-1 = 1\) \\
2 & 1 & \(2\cdot 1+1 = 3\equiv 0\) & \(2\cdot 1-1 = 1\) \\
3 & 2 & \(0\cdot 2+1 = 1\) & \(0\cdot 2-1 = -1\equiv 2\) \\
4 & 1 & \(1\cdot 1+1 = 2\) & \(1\cdot 1-1 = 0\) \\
5 & 2 & \(2\cdot 2+1 = 5\equiv 2\) & \(2\cdot 2-1 = 3\equiv 0\) \\
\bottomrule
\end{tabular}
\end{center}

An equivalent compact form of the table is
\[
\begin{array}{c|cccccc}
n\bmod 6 & 0&1&2&3&4&5\\ \hline
2^n \bmod 3 & 1&2&1&2&1&2\\
\Cul_n \bmod 3 & 1&0&0&1&2&2\\
\Woo_n \bmod 3 & 2&1&1&2&0&0
\end{array}
\]

\section{The two equivalences}

\begin{theorem}[Cullen]\label{thm:cullen}
For every \(n\ge 1\),
\[
3\mid \Cul_n = n2^n+1
\iff
n\equiv 1 \text{ or } 2 \pmod 6 .
\]
Consequently the proportion of Cullen numbers divisible by \(3\) (equivalently,
with least prime factor \(3\)) is exactly
\[
\frac{2}{6}=\frac13 .
\]
\end{theorem}

\begin{proof}
By the residue table, \(\Cul_n \bmod 3 = 0\) exactly in the rows
\(r=n\bmod 6\in\{1,2\}\):
\begin{itemize}
\item \(r=1\): \(\Cul_n\equiv 1\cdot 2+1 = 3\equiv 0\pmod 3\);
\item \(r=2\): \(\Cul_n\equiv 2\cdot 1+1 = 3\equiv 0\pmod 3\);
\end{itemize}
while in the remaining rows
\[
\Cul_n\equiv 1,1,2,2 \pmod 3
\qquad (r=0,3,4,5),
\]
none of which is \(0\). This proves the equivalence. Among the six residue
classes modulo \(6\), exactly two satisfy it, so the proportion is
\(2/6=1/3\).
\end{proof}

\begin{theorem}[Woodall]\label{thm:woodall}
For every \(n\ge 1\),
\[
3\mid \Woo_n = n2^n-1
\iff
n\equiv 4 \text{ or } 5 \pmod 6 .
\]
Consequently the proportion of Woodall numbers divisible by \(3\) (equivalently,
with least prime factor \(3\)) is exactly
\[
\frac{2}{6}=\frac13 .
\]
\end{theorem}

\begin{proof}
By the residue table, \(\Woo_n \bmod 3 = 0\) exactly in the rows
\(r=n\bmod 6\in\{4,5\}\):
\begin{itemize}
\item \(r=4\): \(\Woo_n\equiv 1\cdot 1-1 = 0\pmod 3\);
\item \(r=5\): \(\Woo_n\equiv 2\cdot 2-1 = 3\equiv 0\pmod 3\);
\end{itemize}
while in the remaining rows
\[
\Woo_n\equiv 2,1,1,2 \pmod 3
\qquad (r=0,1,2,3),
\]
none of which is \(0\). This proves the equivalence. Among the six residue
classes modulo \(6\), exactly two satisfy it, so the proportion is
\(2/6=1/3\).
\end{proof}

Observe that the two sets of residues are disjoint and complementary
\emph{within the six residue classes}:
\[
\{1,2\}\ \dot\cup\ \{4,5\} \subsetneq \{0,1,2,3,4,5\},
\]
the two classes \(r=0\) and \(r=3\) being divisible by \(3\) in neither
sequence. This is precisely why the proportions do not complete to \(1\).

\section{The proportions do not sum to \(1\)}

\begin{theorem}[Refutation]\label{thm:refute}
Let
\[
p_{\Cul} = \lim_{N\to\infty}\frac{\#\{1\le n\le N : 3\mid \Cul_n\}}{N},
\qquad
p_{\Woo} = \lim_{N\to\infty}\frac{\#\{1\le n\le N : 3\mid \Woo_n\}}{N}
\]
be the densities of Cullen and Woodall numbers divisible by \(3\). Then
\[
p_{\Cul} = \frac13, \qquad p_{\Woo} = \frac13, \qquad
p_{\Cul}+p_{\Woo} = \frac23 \neq 1 .
\]
Hence the conjectured complementary symmetry modulo \(3\) is false.
\end{theorem}

\begin{proof}
Because divisibility by \(3\) is periodic in \(n\) with period \(6\), the
density of each set is the fraction of residue classes modulo \(6\) it
occupies. By Theorems~\ref{thm:cullen} and \ref{thm:woodall},
\[
p_{\Cul} = \frac{|\{1,2\}|}{6} = \frac{2}{6} = \frac13,
\qquad
p_{\Woo} = \frac{|\{4,5\}|}{6} = \frac{2}{6} = \frac13 .
\]
Therefore
\[
p_{\Cul}+p_{\Woo} = \frac13+\frac13 = \frac23,
\]
and \(\tfrac23\neq 1\). Equivalently, over the common denominator \(6\) the
two numerators \(2,2\) sum to \(4\), and \(4/6 \neq 6/6\) because
\(4\neq 6\). The refutation is unconditional.
\end{proof}

\begin{remark}
The refutation does not depend on the asymptotic count of Cullen primes, nor
on any unproved hypothesis. It is a finite computation on the six residue
classes modulo \(6\). Numerically, over all \(n\le 60000\) one finds densities
\(0.3333\) for each sequence, with the divisible residues exactly \(\{1,2\}\)
for Cullen and \(\{4,5\}\) for Woodall, in agreement with
Theorems~\ref{thm:cullen} and~\ref{thm:woodall}.
\end{remark}

\section{Formal and computational verification}

\paragraph{Lean~4.}
The file \texttt{lean4/Main.lean} formalises the argument in core Lean
(\texttt{import Std}, no Mathlib, no \texttt{sorry}, no \texttt{axiom}, no
\texttt{native\_decide}). It proves the period-\(6\) identities
\begin{align*}
2^n \bmod 3 &= 2^{\,n\bmod 6} \bmod 3,\\
(n2^n+1)\bmod 3 &= ((n\bmod 6)\,2^{\,n\bmod 6}+1)\bmod 3,\\
(n2^n+2)\bmod 3 &= ((n\bmod 6)\,2^{\,n\bmod 6}+2)\bmod 3,
\end{align*}
then the equivalences
\begin{align*}
\texttt{cullen\_div3\_iff} &: (n2^n+1)\bmod 3 = 0 \iff n\bmod 6\in\{1,2\},\\
\texttt{woodall\_div3\_iff} &: (n2^n+2)\bmod 3 = 0 \iff n\bmod 6\in\{4,5\},
\end{align*}
the sequence form \texttt{woodall\_W\_div3\_iff} for
\(\Woo_n=n2^n-1\) with \(n\ge 1\), the parity lemmas
\texttt{cullen\_odd} and \texttt{woodall\_odd}, the residue table
\texttt{residue\_table}, the non-sum statement \texttt{proportions\_sum}
(\(2+2=4\) and \(4\neq 6\) over the common denominator \(6\)), and the
collected theorem \texttt{conjecture\_00000000226\_false}. The axiom audit in
\texttt{lean4/Check.lean} reports only \texttt{propext} and
\texttt{Quot.sound}, with no \texttt{sorryAx} and no
\texttt{Lean.ofReduceBool}.

\paragraph{Python.}
The script \texttt{reproduce.py} (standard library only) verifies both iff
statements for all \(1\le n\le 2000\), prints the residue table, checks the
parity of both sequences, and asserts exactly that
\(\frac13+\frac13=\frac23\neq 1\). It prints \texttt{PASS} and exits with
status \(0\).

\begin{thebibliography}{9}
\bibitem{cun}
J. M. Cullen, \emph{Question 15897}, The Educational Times \textbf{23}
(1905), 534.
\bibitem{woo}
A. J. C. Cunningham and H. J. Woodall, \emph{Factorisation of
\(y^n\mp 1\), \(y=2,3,5,6,7,10,11,12\) up to high powers},
Messenger of Mathematics \textbf{48} (1917), 237--240.
\bibitem{conj}
Conjecture \texttt{00000000226}, \texttt{conjectures/00000000226.md},
this repository.
\end{thebibliography}

\end{document}
